\daytitle{14}{机械波}

\begin{learningbox}{本讲义大致内容}
机械波是“振动在介质中的传播”。今天把 Day 13 的机械振动向空间推广：一个质点振动，会带动周围质点依次振动，于是形成波。重点掌握波长、频率、周期、波速的关系，区分波形图和振动图，并会用微平移法判断质点振动方向。
\end{learningbox}

\section*{一、学习目标}

\begin{enumerate}
  \item 理解机械波的形成条件：振源和介质。
  \item 知道机械波传播的是振动形式和能量，介质质点不随波迁移。
  \item 区分横波和纵波，能识别波峰、波谷、密部、疏部。
  \item 掌握 $T$、$f$、$\lambda$、$v$ 的物理意义和关系 $v=\lambda f$。
  \item 能读懂 $y$-$x$ 波形图和 $y$-$t$ 振动图，并判断质点振动方向。
\end{enumerate}

\section*{二、初高中衔接}

初中见过水波、声波，也知道“声音的传播需要介质”。高中要进一步回答三个问题：

\begin{enumerate}
  \item \textbf{谁在振动？} 介质中的每个质点都在自己的平衡位置附近振动。
  \item \textbf{什么在传播？} 传播的是振动状态和能量，不是介质整体被运走。
  \item \textbf{图像横轴是什么？} $y$-$x$ 图像表示同一时刻不同质点的位置；$y$-$t$ 图像表示同一质点随时间的振动。
\end{enumerate}

这节课最容易混淆的地方，正是“波向前走”和“质点只在附近振动”这两件事要同时成立。

\section*{三、核心概念}

\begin{conceptbox}{机械波的形成}

\textbf{定义：} 机械振动在介质中由近及远传播，形成机械波。

\textbf{形成条件：}
\begin{itemize}
  \item \textbf{振源：} 先有一个持续或短暂的振动。
  \item \textbf{介质：} 相邻质点之间能相互作用，把振动传下去。
\end{itemize}

\begin{center}
% TikZ 图：横波的形成
% Day 14：机械波
\begin{tikzpicture}[x=0.9cm,y=0.75cm,>=Latex]
  \draw[->] (-0.3,0) -- (7.0,0) node[right] {$x$};
  \draw[dashed,gray] (0,0) -- (6.4,0);
  \draw[very thick,red,smooth,domain=0:6.28,samples=90]
    plot (\x,{sin(\x r)});

  \foreach \x in {0,0.52,...,6.24} {
    \pgfmathsetmacro{\y}{sin(\x r)}
    \draw[gray!55] (\x,0) -- (\x,\y);
    \filldraw[black] (\x,\y) circle (1.4pt);
  }

  \draw[<->,thick,blue] (1.57,-1.28) -- (1.57,1.28);
  \node[blue,above] at (1.57,1.28) {质点振动};
  \draw[->,thick] (5.2,1.35) -- (6.35,1.35) node[right] {波速 $v$};
\end{tikzpicture}


\figcaption{横波的形成：介质质点上下振动，波形向右传播}\label{fig:wave-transverse-formation}
\end{center}

\textbf{关键认识：} 质点并不跟着波一直向前跑。波向右传播时，每个质点只在自己的平衡位置附近振动。

\end{conceptbox}

\begin{conceptbox}{横波与纵波}

\textbf{横波：} 质点振动方向与波的传播方向垂直。绳波、水面波的直观图像常画成横波。

\textbf{纵波：} 质点振动方向与波的传播方向在同一直线上。声波是典型纵波，空气中出现疏部和密部。

\begin{center}
% TikZ 图：纵波的疏密
% Day 14：机械波
\begin{tikzpicture}[x=0.9cm,y=0.85cm,>=Latex]
  \draw[->] (-0.3,0) -- (7.4,0) node[right] {$x$};
  \draw[->,thick] (5.9,1.05) -- (7.0,1.05) node[right] {波速 $v$};

  \foreach \x in {0.15,0.55,1.10,1.42,1.70,1.95,2.20,2.48,3.05,3.70,4.18,4.46,4.72,5.00,5.28,5.86,6.45} {
    \filldraw[black] (\x,0) circle (1.8pt);
  }

  \draw[<->,blue,thick] (0.45,-0.55) -- (1.05,-0.55);
  \draw[<->,blue,thick] (3.35,-0.55) -- (3.95,-0.55);
  \node[blue,below] at (2.2,-0.55) {质点沿传播方向往复振动};

  \draw[|-|] (1.25,0.48) -- (2.55,0.48);
  \node[above] at (1.9,0.65) {密部};
  \draw[|-|] (2.85,0.48) -- (3.95,0.48);
  \node[above] at (3.4,0.65) {疏部};
  \draw[|-|] (4.10,0.48) -- (5.35,0.48);
  \node[above] at (4.75,0.65) {密部};
\end{tikzpicture}


\figcaption{纵波：疏部和密部沿传播方向前进}\label{fig:wave-longitudinal-density}
\end{center}

\textbf{判断方法：} 只比较“质点振动方向”和“波传播方向”。不要只看图像像不像正弦曲线。

\end{conceptbox}

\begin{conceptbox}{描述机械波的物理量}

\begin{center}
% TikZ 图：波长、振幅和波速
% Day 14：机械波
\begin{tikzpicture}[x=0.95cm,y=0.82cm,>=Latex]
  \draw[->] (-0.35,0) -- (6.8,0) node[right] {$x$};
  \draw[->] (0,-1.45) -- (0,1.65) node[above] {$y$};
  \node[below left] at (0,0) {$O$};

  \draw[very thick,red,smooth,domain=0:6.2,samples=100]
    plot (\x,{sin(120*\x)});

  \draw[dashed,gray] (0.75,0) -- (0.75,1);
  \draw[dashed,gray] (3.75,0) -- (3.75,1);
  \draw[<->] (0.75,1.25) -- (3.75,1.25) node[midway,above] {$\lambda$};

  \draw[<->] (0.25,0) -- (0.25,1) node[midway,right] {$A$};
  \node[above] at (0.75,1.05) {波峰};
  \node[below] at (2.25,-1.05) {波谷};
  \draw[->,thick] (5.0,1.25) -- (6.15,1.25) node[right] {$v$};
\end{tikzpicture}


\figcaption{波的基本物理量：振幅 $A$、波长 $\lambda$、波速 $v$}\label{fig:wave-parameters}
\end{center}

\begin{itemize}
  \item \textbf{振幅 $A$：} 介质质点离开平衡位置的最大距离，反映振动强弱。
  \item \textbf{周期 $T$：} 一个质点完成一次全振动的时间。
  \item \textbf{频率 $f$：} 一个质点单位时间完成全振动的次数，$f=\dfrac{1}{T}$。
  \item \textbf{波长 $\lambda$：} 沿传播方向，相邻两个同相质点之间的距离。
  \item \textbf{波速 $v$：} 波形向前传播的速度。
\end{itemize}

\textbf{核心公式：}
\[
\boxed{v=\frac{\lambda}{T}=\lambda f}
\]

\end{conceptbox}

\begin{conceptbox}{波形图与振动图}

\begin{center}
% TikZ 图：波形图和振动图对比
% Day 14：机械波
\begin{tikzpicture}[x=0.82cm,y=0.72cm,>=Latex]
  % 左：波形图
  \begin{scope}
    \draw[->] (-0.25,0) -- (3.6,0) node[right] {$x$};
    \draw[->] (0,-1.25) -- (0,1.35) node[above] {$y$};
    \draw[very thick,red,smooth,domain=0:3.2,samples=70]
      plot (\x,{sin(180*\x)});
    \draw[<->] (0.5,1.15) -- (2.5,1.15) node[midway,above] {$\lambda$};
    \node[font=\footnotesize,align=center,text width=3.6cm] at (1.65,-1.55) {某一时刻\\整条介质的位置};
    \node[font=\bfseries] at (1.65,1.75) {$y$-$x$ 波形图};
  \end{scope}

  % 右：振动图
  \begin{scope}[shift={(5.2,0)}]
    \draw[->] (-0.25,0) -- (3.6,0) node[right] {$t$};
    \draw[->] (0,-1.25) -- (0,1.35) node[above] {$y$};
    \draw[very thick,blue,smooth,domain=0:3.2,samples=70]
      plot (\x,{sin(180*\x)});
    \draw[<->] (0.5,1.15) -- (2.5,1.15) node[midway,above] {$T$};
    \node[font=\footnotesize,align=center,text width=3.6cm] at (1.65,-1.55) {某一质点\\随时间的振动};
    \node[font=\bfseries] at (1.65,1.75) {$y$-$t$ 振动图};
  \end{scope}
\end{tikzpicture}


\figcaption{$y$-$x$ 波形图与 $y$-$t$ 振动图的区别}\label{fig:wave-yx-yt-compare}
\end{center}

\begin{itemize}
  \item $y$-$x$ 波形图：固定一个时刻，看整条介质上各质点的位置。横轴是位置 $x$，可以读出波长 $\lambda$。
  \item $y$-$t$ 振动图：固定一个质点，看它随时间怎样振动。横轴是时间 $t$，可以读出周期 $T$。
\end{itemize}

\textbf{一句话记忆：} 波形图看“一排人此刻在哪里”，振动图看“一个人随时间怎么动”。

\end{conceptbox}

\begin{conceptbox}{传播方向与质点振动方向}

在波形图上，波的传播方向和某个质点的振动方向不是同一个方向。判断时常用“微平移法”：

\begin{center}
% TikZ 图：微平移法判断质点振动方向
% Day 14：机械波
\begin{tikzpicture}[x=1.0cm,y=0.85cm,>=Latex]
  \draw[->] (-0.3,0) -- (6.4,0) node[right] {$x$};
  \draw[->] (0,-1.35) -- (0,1.45) node[above] {$y$};
  \draw[very thick,red,smooth,domain=0:6.0,samples=100]
    plot (\x,{sin(120*\x)});
  \draw[dashed,gray,smooth,domain=0.25:6.25,samples=100]
    plot (\x,{sin(120*(\x-0.25))});

  \draw[->,thick] (4.7,1.25) -- (5.8,1.25) node[right] {波向右传播};
  \coordinate (P) at (1.40,{sin(120*1.40)});
  \coordinate (Q) at (1.40,{sin(120*(1.40-0.25))});
  \filldraw[black] (P) circle (1.8pt) node[below left] {$P$};
  \draw[->,thick,blue] (P) -- (Q) node[midway,right] {$v_P$};
  \node[gray] at (3.25,-1.15) {虚线：极短时间后的波形};
\end{tikzpicture}


\figcaption{微平移法：把波形沿传播方向平移一点，看质点下一刻往哪里走}\label{fig:wave-direction-shift}
\end{center}

\begin{enumerate}
  \item 先确定波向哪边传播。
  \item 把整条波形沿传播方向平移很小一段。
  \item 观察原来位置的质点要到达新波形上的对应高度，是向上还是向下。
\end{enumerate}

\textbf{结论：} 对向右传播的横波，斜率为负处质点向上运动，斜率为正处质点向下运动；向左传播时相反。

\end{conceptbox}

\begin{conceptbox}{同相点与反相点}

同相表示两个质点的振动步调完全一致：同时到达波峰、同时经过平衡位置且运动方向相同。

\begin{center}
% TikZ 图：同相点与反相点
% Day 14：机械波
\begin{tikzpicture}[x=0.95cm,y=0.82cm,>=Latex]
  \draw[->] (-0.3,0) -- (6.6,0) node[right] {$x$};
  \draw[->] (0,-1.35) -- (0,1.45) node[above] {$y$};
  \draw[very thick,red,smooth,domain=0:6.15,samples=100]
    plot (\x,{sin(120*\x)});

  \foreach \x/\name/\pos in {0.75/A/above,2.25/C/below,3.75/B/above} {
    \filldraw[black] (\x,{sin(120*\x)}) circle (2pt) node[\pos] {$\name$};
  }

  \draw[<->] (0.75,1.2) -- (3.75,1.2) node[midway,above] {$\lambda$};
  \draw[<->] (0.75,-1.2) -- (2.25,-1.2) node[midway,below] {$\lambda/2$};
  \node[anchor=west] at (4.35,0.8) {$A$ 与 $B$ 同相};
  \node[anchor=west] at (4.35,-0.8) {$A$ 与 $C$ 反相};
\end{tikzpicture}


\figcaption{相距 $\lambda$ 的质点同相，相距 $\lambda/2$ 的质点反相}\label{fig:wave-phase-points}
\end{center}

\begin{itemize}
  \item 相距 $n\lambda$ 的两个质点同相，$n=0,1,2,\cdots$。
  \item 相距 $\dfrac{2n+1}{2}\lambda$ 的两个质点反相。
\end{itemize}

\end{conceptbox}

\section*{四、公式与适用条件}

\begin{center}
\begin{tabular}{lll}
\toprule
公式或关系 & 物理意义 & 使用提醒 \\
\midrule
$f=\dfrac{1}{T}$ & 频率与周期互为倒数 & 同一个质点的振动周期 \\
$v=\dfrac{\lambda}{T}=\lambda f$ & 波速、波长、频率关系 & 机械波普遍适用 \\
$\Delta x=v\Delta t$ & 波形在时间 $\Delta t$ 内前进的距离 & $v$ 是波速，不是质点速度 \\
$\Delta x=n\lambda$ & 两质点同相 & $n$ 为整数 \\
$\Delta x=\dfrac{2n+1}{2}\lambda$ & 两质点反相 & $n$ 为整数 \\
\bottomrule
\end{tabular}
\end{center}

\textbf{特别提醒：} 在同一种介质中，波速主要由介质性质决定。振源频率改变时，波长随之改变，使 $v=\lambda f$ 仍成立。

\section*{五、典型模型}

\subsection*{模型 1：绳上的横波}

绳上每个质点上下振动，波形沿绳传播。解题时先读图得到 $A$、$\lambda$，再用 $v=\lambda f$ 或 $\Delta x=v\Delta t$ 处理传播。

\subsection*{模型 2：声波}

声波在空气中是纵波。空气分子沿传播方向往复振动，疏部和密部向前传播。声波可以传递能量，但空气整体不会随声音从声源跑到人耳。

\subsection*{模型 3：波形图方向判断}

题目常给某时刻波形，并告诉波向左或向右传播。判断某点下一刻向上还是向下时，用微平移法，不要凭“离波峰近”猜。

\subsection*{模型 4：相位关系}

相邻波峰、相邻波谷、一个波峰到下一个波峰的距离都是 $\lambda$；一个波峰到相邻波谷的距离是 $\lambda/2$，这两点反相。

\section*{六、例题精讲}

\begin{examplebox}{波长、周期、波速的计算}

\textbf{题意分析：} 一列横波在绳上向右传播。某时刻波形图读出振幅 $A=4\ \mathrm{cm}$，波长 $\lambda=0.80\ \mathrm{m}$。若振源频率 $f=5.0\ \mathrm{Hz}$，求周期、波速，以及 $2.0\ \mathrm{s}$ 内波形传播的距离。

\textbf{物理模型：} 机械横波的基本物理量计算。

\textbf{解题过程：}
\[
T=\frac{1}{f}=\frac{1}{5.0}=0.20\ \mathrm{s}
\]
\[
v=\lambda f=0.80\times 5.0=4.0\ \mathrm{m/s}
\]
\[
\Delta x=v\Delta t=4.0\times 2.0=8.0\ \mathrm{m}
\]

\textbf{方法总结：} $A$ 描述质点振动强弱，不参与波速计算；波速由 $\lambda$ 和 $f$ 的乘积给出。

\end{examplebox}

\begin{examplebox}{波形平移求下一时刻位置}

\textbf{题意分析：} 一列波沿 $x$ 轴正方向传播，波速 $v=2.0\ \mathrm{m/s}$，波长 $\lambda=4.0\ \mathrm{m}$。$t=0$ 时，某波峰在 $x=1.0\ \mathrm{m}$ 处。求 $0.50\ \mathrm{s}$ 后这个波峰的位置。

\textbf{物理模型：} 波形整体沿传播方向平移。

\textbf{解题过程：}

波形在 $0.50\ \mathrm{s}$ 内前进：
\[
\Delta x=v\Delta t=2.0\times 0.50=1.0\ \mathrm{m}
\]

波向右传播，所以波峰位置变为：
\[
x'=1.0+1.0=2.0\ \mathrm{m}
\]

\textbf{方法总结：} 波峰、波谷不是固定质点，而是波形上的特征位置。波向右传播时，整个波形向右平移。

\end{examplebox}

\begin{examplebox}{由波形图判断质点运动方向}

\textbf{题意分析：} 某时刻一列横波沿 $x$ 轴正方向传播。图上某点 $P$ 位于平衡位置上方，且所在处波形斜率为负。判断 $P$ 点此刻向上还是向下运动。

\textbf{物理模型：} 右行横波的微平移判断。

\textbf{解题过程：}

波向右传播，把波形向右平移一小段。原来 $P$ 点所在的 $x$ 位置，将对应到平移后波形上更高的位置，因此 $P$ 点要向上运动。

\[
\boxed{P\text{ 点向上运动}}
\]

\textbf{方法总结：} 对右行波，斜率为负处质点向上，斜率为正处质点向下。对左行波，判断结果相反。

\end{examplebox}

\section*{七、变式训练}

\begin{enumerate}
  \item 一列机械波的频率 $f=20\ \mathrm{Hz}$，波长 $\lambda=0.50\ \mathrm{m}$。求波速和周期。

  \item 一列波的波速为 $340\ \mathrm{m/s}$，频率为 $170\ \mathrm{Hz}$。求波长，并说明这可能对应哪类生活中的机械波。

  \item 一列横波向右传播，某时刻点 $A$ 处波形斜率为正，点 $B$ 处波形斜率为负。判断 $A$、$B$ 此刻的振动方向。

  \item 一列简谐波中，$P$、$Q$ 两质点相距 $1.5\ \mathrm{m}$。若波长 $\lambda=1.0\ \mathrm{m}$，判断 $P$、$Q$ 是同相还是反相；若 $\lambda=3.0\ \mathrm{m}$ 呢？
\end{enumerate}

\section*{八、易错点}

\begin{mistakebox}
\begin{enumerate}
  \item \textbf{认为波把介质带走：} 机械波传播的是振动状态和能量，介质质点只在平衡位置附近振动。

  \item \textbf{混淆波速和质点速度：} 波速 $v$ 是波形传播速度；质点速度是某个介质质点振动时的瞬时速度，大小和方向都随时间变化。

  \item \textbf{混淆 $y$-$x$ 图像和 $y$-$t$ 图像：} 前者横轴是位置，读波长；后者横轴是时间，读周期。

  \item \textbf{误以为 $v=\lambda f$ 中频率决定波速：} 在同一介质中，波速主要由介质决定。频率变大时，通常是波长变小。

  \item \textbf{把相邻零点距离当作波长：} 相邻两个同向过平衡位置的点相距 $\lambda$；相邻两个零点通常相距 $\lambda/2$。

  \item \textbf{横波纵波按图形形状判断：} 横波和纵波看的是质点振动方向与传播方向的关系，不是看画出来像不像“横着的曲线”。
\end{enumerate}
\end{mistakebox}

\clearpage
\section*{九、竞赛入门补充}

\begin{competitionbox}
\textbf{波的叠加原理}

当两列波在同一介质中相遇时，每个质点的实际位移等于两列波分别引起位移的代数和。相遇后，两列波仍按原来的方向传播。

\begin{center}
% TikZ 图：两列波的叠加
% Day 14：机械波
\begin{tikzpicture}[x=0.8cm,y=0.75cm,>=Latex]
  % 第一行：相遇前
  \draw[->] (-0.2,0) -- (8.2,0) node[right] {$x$};
  \draw[very thick,red] (0.7,0) -- (1.2,0.9) -- (1.7,0);
  \draw[very thick,blue] (5.9,0) -- (6.4,0.9) -- (6.9,0);
  \draw[->,red,thick] (1.9,0.55) -- (2.8,0.55);
  \draw[->,blue,thick] (5.7,0.55) -- (4.8,0.55);
  \node[left] at (-0.1,0.7) {相遇前};

  % 第二行：重叠时
  \begin{scope}[shift={(0,-1.75)}]
    \draw[->] (-0.2,0) -- (8.2,0) node[right] {$x$};
    \draw[dashed,red] (3.5,0) -- (4.0,0.9) -- (4.5,0);
    \draw[dashed,blue] (3.5,0) -- (4.0,0.9) -- (4.5,0);
    \draw[very thick,purple] (3.5,0) -- (4.0,1.8) -- (4.5,0);
    \node[left] at (-0.1,0.8) {叠加时};
    \node[right] at (4.7,1.25) {$y=y_1+y_2$};
  \end{scope}

  % 第三行：相遇后
  \begin{scope}[shift={(0,-3.5)}]
    \draw[->] (-0.2,0) -- (8.2,0) node[right] {$x$};
    \draw[very thick,blue] (0.7,0) -- (1.2,0.9) -- (1.7,0);
    \draw[very thick,red] (5.9,0) -- (6.4,0.9) -- (6.9,0);
    \draw[->,blue,thick] (1.9,0.55) -- (1.0,0.55);
    \draw[->,red,thick] (5.7,0.55) -- (6.6,0.55);
    \node[left] at (-0.1,0.7) {相遇后};
  \end{scope}
\end{tikzpicture}


\figcaption{两列向上位移脉冲相遇时，位移相加；相遇后继续传播}\label{fig:wave-superposition}
\end{center}

\textbf{模型意识：} 叠加不是“碰撞后合成一列波”，而是同一质点同时参与两个振动。若两列周期相同、方向相反的波长期叠加，就会出现驻波的雏形：有些点几乎不动，叫波节；有些点振幅最大，叫波腹。

\textbf{竞赛小贴士：} 遇到两列波叠加，先分别画出每列波在该时刻造成的位移，再按正负相加，不要把波峰、波谷当作物体相撞。
\end{competitionbox}

\section*{十、今日小结}

\begin{enumerate}
  \item 机械波需要振源和介质，传播的是振动形式和能量。
  \item 横波中质点振动方向与传播方向垂直；纵波中二者在同一直线上。
  \item 波的基本关系为 $v=\dfrac{\lambda}{T}=\lambda f$。
  \item $y$-$x$ 波形图读 $\lambda$，$y$-$t$ 振动图读 $T$。
  \item 判断质点振动方向时，用微平移法；判断相位时，看两点间距与 $\lambda$ 的关系。
\end{enumerate}
